\begin{problem}[Warmup]{(a)}
    Give examples of the following, or explain why such an example cannot exist. Be sure to specify explicitly the vector space, the field of scalars, and the vectors required:
    \begin{enumerate}
        \item[(a)] A spanning set that isn't a basis.
        \item[(b)] A basis that isn't a spanning set.
        \item[(c)] Two distinct bases for the same vector space.
        \item[(d)] Two distinct bases for the same vector space, one of which is a proper subset of the other.
        \item[(e)] A set of three linearly dependent vectors, any two of which are independent.
        \item[(f)] A set of four linearly dependent vectors, any three of which are independent.
    \end{enumerate}
\end{problem}
\unsolved

\begin{problem}[Warmup]{(b)}
    Why could the zero vector never belong to a basis (for a vector space with more than one element)?
\end{problem}
\unsolved

\begin{problem}[Warmup]{(c)}
    What is the dimension of $\mathbb{C}$ as a vector space over $\mathbb{R}$? What is the dimension of $\mathbb{C}$ as a vector space over $\mathbb{C}$?
\end{problem}
\unsolved

\begin{problem}{1}
    Prove that $\sqrt{2}$ and $\sqrt{3}$ are linearly independent elements of the vector
    space $\mathbb{R}$, over $\mathbb{Q}$.
\end{problem}
\unsolved

\begin{problem}{2}
    (Here, we extend Exercise 1.) Prove that $\{1, \sqrt{2}, \sqrt{3}\}$ is an independent
    subset of the vector space $\mathbb{R}$ over $\mathbb{Q}$. Is it a basis for
    $\mathbb{R}$?
\end{problem}
\unsolved

\begin{problem}{3}
    Consider the set $V$, consisting of all polynomials of degree 0 or 1, with coefficients
    from $\mathbb{C}$. Prove that this is both a vector space over $\mathbb{C}$, and over
    $\mathbb{R}$. Find a basis in each case. Does its dimension stay the same, when we
    change the scalar field?
\end{problem}
\unsolved

\begin{problem}{4}
    Find a basis for the vector space given in Exercise 40.7.
\end{problem}
\unsolved

\begin{problem}{5}
    Find a basis for the vector space given in Exercise 40.8.
\end{problem}
\unsolved

\begin{problem}{6}
    Give a basis for the vector space given in Exercise 40.10.
\end{problem}
\unsolved

\begin{problem}{7}
    Show that $\mathbb{Q}[x]$ over $\mathbb{Q}$ is not finite dimensional by showing that no
    finite subset of $\mathbb{Q}[x]$ is a spanning set.
\end{problem}
\unsolved

\begin{problem}{8}
    Prove Theorem 41.2.
\end{problem}
\unsolved

\begin{problem}{9}
    \begin{enumerate}
        \item[(a)] Define a \textbf{subspace} of a vector space.
        \item[(b)] Let $\mathcal{S}$ be a subset of the vector space $V$. Show that the set
        of linear combinations of vectors from $\mathcal{S}$ is a subspace of $V$. (Note
        that this subspace is all of $V$ if and only if $\mathcal{S}$ spans $V$.)
    \end{enumerate}
\end{problem}
\unsolved

\begin{problem}{10}
    Determine which of the following are subspaces of the vector space $M_2(\mathbb{R})$
    discussed in Exercise 40.7.
    \begin{enumerate}
        \item[(a)] Matrices with determinant 1.
        \item[(b)] Matrices with determinant 0.
        \item[(c)] The diagonal matrices.
        \item[(d)] Matrices with integer entries.
        \item[(e)] Matrices of the form
        $\begin{pmatrix} a & a + b \\ b & 0 \end{pmatrix}$.
        \item[(f)] Matrices with first row, first column entry of 0.
    \end{enumerate}
\end{problem}
\unsolved

\begin{problem}{11}
    Determine which of the following are subspaces of the vector space of all real-valued
    functions discussed in Exercise 40.11.
    \begin{enumerate}
        \item[(a)] All differentiable functions.
        \item[(b)] All functions $f$ with $f(2) = 0$.
        \item[(c)] All linear functions. (That is, functions of the form $f(x) = a + bx$.)
        \item[(d)] All polynomial functions.
    \end{enumerate}
\end{problem}
\unsolved

\begin{problem}{12}
    Check Example 41.12: that is, prove that $\{1 + i, \sqrt{2} + 4i\}$ is a spanning set
    for $\mathbb{C}$ over $\mathbb{R}$.
\end{problem}
\unsolved

\begin{problem}{13}
    Find a subset of the set
    \[\{(1, 1, 0), (2, 3, -1), (0, 1, -1), (1, 4, 2), (0, 0, 3)\}\]
    that is a basis for $\mathbb{R}^3$ using the technique described in this chapter.
\end{problem}
\unsolved

\begin{problem}{14}
    Find a subset of the set
    \[\{(-1, 0, -1), (0, 0, 1), (1, 0, 0), (2, 3, 0), (0, 1, 0)\}\]
    that is a basis for $\mathbb{R}^3$ using the technique in this chapter; this is the same
    set we used in Section 41.3 to demonstrate this technique. Take the vectors in the
    order given here, to see explicitly that we obtain a different subset than in the text.
\end{problem}
\unsolved

\begin{problem}{15}
    Let $V$ be a vector space. Show that if $\mathbf{v}$ is a linear combination of the
    vectors $\mathbf{v}_1, \mathbf{v}_2, \ldots, \mathbf{v}_n \in V$ and each
    $\mathbf{v}_i$ is in turn a linear combination of
    $\mathbf{w}_1, \mathbf{w}_2, \ldots, \mathbf{w}_m$, then $\mathbf{v}$ is a linear
    combination of $\mathbf{w}_1, \mathbf{w}_2, \ldots, \mathbf{w}_m$.
\end{problem}
\unsolved

\begin{problem}{16}
    Prove that a basis for a vector space $V$ is a \emph{maximal independent set}: That
    is, it is a set $\mathcal{B} \subseteq V$ of independent vectors, so that
    $\mathcal{B} \cup \{w\}$ is a dependent set, whenever $w \in V \setminus \mathcal{B}$.
\end{problem}
\unsolved

\begin{problem}{17}
    Prove that a maximal independent set of vectors of a vector space is necessarily a
    basis (that is, prove the converse of Exercise 16).
\end{problem}
\unsolved

\begin{problem}{18}
    Use the technique in the proof of Corollary 41.5 to build a basis for $\mathbb{R}^4$
    that contains the linearly independent set
    \[\{(1, 0, -1, 0), (-2, 1, 2, 0)\}.\]
\end{problem}
\unsolved
